\paragraph{An actual completed projector lift of the four-dimensional source module.}
The preceding semisimple obstruction does not prevent the following
specific construction with infinite tangle complexes.  This construction
uses the actual source representation and its generators.  Its
decategorification is a completed Euler group, with the convention stated
below, rather than an ordinary split Grothendieck group of a semisimple
category.

Write $V_j$ for the $(j+1)$-dimensional type-one representation of
$U_q\mathfrak{sl}_2$ over $F_0=\mathbb C(q)$.  Put
\[
 \mathcal I=\operatorname{Hom}_{U_q\mathfrak{sl}_2}
       (V_3^{\otimes4},F_0),\qquad
 \dim_{F_0}\mathcal I=4.
                                                               \tag{P-F52bs}
\]
The four fusion vectors $\psi_j$, $j=0,2,4,6$, join the first
two and last two factors through the common channel $V_j$.
We fix any nonzero standard trivalent normalization for each vector.
The normalized source operator $Q_i/f$ is exactly the sum of
the channel-$0$ and channel-$4$ projections on factors $i,i+1$,
for $i=1,2$.

The primary categorical inputs are Cooper--Krushkal,
\emph{Categorification of the Jones--Wenzl projectors},
Definitions~2.8--2.9 and Lemma~2.10, printed p.~9,
Theorem~3.2, p.~12, and Theorem~6.2, p.~21:
\url{https://arxiv.org/pdf/1005.5117}.
Their completed Euler convention identifies a complex with its
degreewise alternating sum in
$\widehat R=\mathbb Z[[q]][q^{-1}]$.
Precisely, it is the $q$-monotone subcategory
$\operatorname{Kom}_0$ in their Definition~2.8, with the
convergent graded Euler sums and infinite Euler relations
of Definition~2.9.  The complexes are bounded below in
homological degree; this convention does not replace them
by arbitrary unbounded complexes.
Their category $\operatorname{Kom}(3,3,3,3)$ uses four boundary
Jones--Wenzl projector complexes $\mathbf P_3$.
We use the corresponding homotopy category and its Euler group,
denoted $K^{\mathrm{Eul}}$ here.
Thus
\[
 K^{\mathrm{Eul}}(\operatorname{Kom}(3,3,3,3))
       \otimes_{\widehat R}\mathbb C((q))
       \simeq\mathcal I\otimes_{F_0}\mathbb C((q)).
                                                               \tag{P-F52bt}
\]
The rational subspace on the right is the $F_0$-span of the four
fusion vectors.  All the operators below preserve that subspace.
No assertion that an unrestricted unbounded homotopy category has
this ordinary triangulated Grothendieck group is intended.
In particular (P-F52bt) is the isomorphism supplied by the cited
completed Euler convention, not merely a surjective character map.
The inclusion $F_0\hookrightarrow\mathbb C((q))$ is the Laurent
expansion at zero.  The four rational fusion classes are linearly
independent after this extension, and their $F_0$-span is the
specified four-dimensional comparison space.

Cooper--Hogancamp, \emph{An exceptional collection for Khovanov
homology}, Definition~8.5 and Theorem~8.7, printed p.~39,
and Propositions~8.13--8.14,
p.~41, construct the higher projector complexes
$\mathbf P_{n,j}$, prove their orthogonal idempotence, and permit
them to pass across any compatible tangle complex:
\url{https://arxiv.org/pdf/1209.1002}.
Let $\mathscr E_i$ be the functor obtained by inserting
\[
                  \mathbf P_{6,0}\oplus\mathbf P_{6,4}
                                                               \tag{P-F52bu}
\]
on the six strands belonging to boundary bundles $i,i+1$,
with boundary $\mathbf P_3$'s retained.  Equivalently, sandwich
this insertion between the boundary projector complexes.
Let $\mathcal H_{Ff}$ be the finite direct sum of grading shifts
whose character is the nonnegative Laurent polynomial $Ff$, and put
\[
             \mathscr Q_i=\mathcal H_{Ff}\otimes\mathscr E_i.
                                                               \tag{P-F52bv}
\]
These are actual additive functors on the completed tangle category.
They admit natural homotopy equivalences
$\mathscr E_i^2\simeq\mathscr E_i$ and
$\mathscr Q_i^2\simeq\mathcal H_{Ff}\otimes\mathscr Q_i$.
Their actions on the rational Euler subspace are exactly $FQ_i$.

There is an explicit $F_0$-linear isomorphism from the actual module
$V=F_0\mathcal V$ of (P-F52ax) to that rational subspace:
\[
                \Phi(v)=v\psi_0,
 \qquad \Phi(FQ_i v)=[\mathscr Q_i]\Phi(v).
                                                               \tag{P-F52bw}
\]
In this formula the convenient word acts by the source tensor
operators; the same formula, expanded in the fusion basis, specifies
every column of the change of basis.  If invariants are written as
maps out of the tensor product, use the word-reversing
anti-involution fixing $Q_1,Q_2$ to turn precomposition into a
left action.  This convention gives the displayed generator
intertwining and respects the order of every word.

\emph{Proof of the comparison and functor assertions.}
The Clebsch--Gordan decomposition is
$V_3\otimes V_3=V_6\oplus V_4\oplus V_2\oplus V_0$.
The source's four braid eigenvalues, in this order, are
$q^{9/2},-q^{-3/2},q^{-11/2},-q^{-15/2}$.
The signs therefore identify its negative spectral projector with
the two channels in (P-F52bu).  This is the original definition
and normalization in GCT VII, \S7.1, printed pp.~17--18,
equations~(16)--(18); it does not use the quarantined
two-by-two display on p.~20.
Each common channel in the two pairings contributes one invariant,
proving the dimension and the fusion basis in (P-F52bs).

For completeness, the particular source module on this invariant
space is identified in GCT VII, \S7.1.3, equation~(25), p.~21,
and equation~(26), p.~22: the two conjugate source blocks
$T_1,T_2$ are precisely the multiplicity blocks for the
$GL_2$ highest weight $(6,3)$ in $X^{\otimes3}$.
Restriction to $\mathfrak{sl}_2$ makes this the $V_3$ multiplicity
space, which is dual to (P-F52bs).  The word-reversal convention
just specified identifies the relevant left actions.  Alternatively,
the nondegenerate symmetric form (P-F52bi), for which both
generators are symmetric, identifies the module and its
word-reversal contragredient over $F_0$.
Consequently the central operator $G$ on $\mathcal I$ has
minimal polynomial $X^2-LX+N_0$, with conjugate values
$g_1,g_2$, each twice.  The same statement can be checked directly
on the kernel of the raising operator in tensor weight $3$ using
the original Figure~1 matrix; the exact finite certificate is
\texttt{check\_nonstandard\_projector\_intertwiner.py}.

The source module $V$ is simple over $F_0$.  Here is a direct
check needed for the map (P-F52bw).  Its center contains the
quadratic field $K=F_0(g_1)$, and $V$ is two-dimensional over $K$.
In the basis consisting of a nonzero $Q_1$-image vector and its
$Q_2$-image, the two unscaled generators are
\[
        Q_1=\begin{pmatrix}f&g_1\\0&0\end{pmatrix},\qquad
        Q_2=\begin{pmatrix}0&0\\1&f\end{pmatrix}.
\]
A common invariant line would have to be an image or kernel
line for both rank-one operators.  Their four possibilities
give $g_1=0$ or $g_1=f^2$, except for the distinct image lines
which never agree.  Neither equality holds, since $g_1\notin F_0$.
Thus there is no proper $K$-submodule, and any $F_0$-submodule
would be $K$-stable by the central action.

The vector $\psi_0$ satisfies $Q_1\psi_0=f\psi_0$.
The actual first lattice vector is
$v_0=(2G-L)h_3(G)Q_1$.
Both $2G-L$ and $h_3(G)$ are invertible on $\mathcal I$:
their values at $g_1,g_2$ are nonzero, because $D\ne0$ and
$g_1,g_2\notin F_0$, whereas $g_3,g_4,g_5\in F_0$.
Hence $v_0\psi_0\ne0$.
More precisely, both conjugate central components are nonzero.
If one of
$(G-g_2)\psi_0/(g_1-g_2)$ and
$(G-g_1)\psi_0/(g_2-g_1)$ vanished after scalar extension to $K$,
the nonzero $F_0$-vector $\psi_0$ would be a $G$-eigenvector
with eigenvalue $g_2$ or $g_1$.  Any nonzero rational coordinate
would then express that eigenvalue as an element of $F_0$,
a contradiction.  The invertible central factors in $v_0$
preserve this nonvanishing in both source blocks.
Evaluation $v\mapsto v\psi_0$ is a nonzero module map from
the simple four-dimensional $V$ to the four-dimensional
$\mathcal I$, and is therefore an isomorphism.  This proves
(P-F52bw) without selecting a purely numerical change of basis.

At the categorical level, Proposition~8.14 applies to the fixed
six-strand complex $\mathbf P_3\sqcup\mathbf P_3$.
It lets each higher channel projector pass through the two
boundary projectors.  In combination with boundary idempotence
and Proposition~8.13, it proves that the sandwiched complex
in (P-F52bu) squares to itself up to homotopy.
Choose such a homotopy equivalence of the fixed inserted
tangle complexes.  Gluing it to a variable input complex
gives a natural homotopy equivalence of functors.
The finite direct sums of shifts in (P-F52bv) preserve all
these constructions.  Tensor products and finite sums of
the bounded-below, convergent complexes used here retain
well-defined completed Euler classes.  Those classes multiply
under gluing, so (P-F52bu) acts by $Q_i/f$ and
(P-F52bv) acts by $FQ_i$, as asserted.
\hfill$\square$

\paragraph{What completion and Euler decategorification leave open.}
The required completion is already visible in the boundary
projector, before either nonstandard generator is inserted.
The Jones--Wenzl recurrence of Cooper--Krushkal, \S2.2, p.~5,
gives
\[
 p_3=1-\frac{[2]}{[3]}(e_1+e_2)
              +\frac1{[3]}(e_1e_2+e_2e_1),\qquad
 [3]=q^2+1+q^{-2}.
                                                               \tag{P-F52bx}
\]
The five displayed diagrams are the Temperley--Lieb basis.
The coefficient $1/[3]$ has a pole at a nonzero root of
$q^4+q^2+1$, so it is not a Laurent polynomial.
A bounded complex of finite sums of shifted ordinary diagrams
has Laurent-polynomial Euler coefficients.  It therefore
cannot represent this boundary projector.
Moreover its expansion at $q=0$
\[
       \frac1{[3]}=\frac{q^2(1-q^2)}{1-q^6}
          =\sum_{r\ge0}(q^{6r+2}-q^{6r+4})
                                                               \tag{P-F52by}
\]
is sent by termwise grading reversal to a series unbounded
in the negative direction, outside $\mathbb C((q))$.
This distinguishes bar invariance of the rational function
from preservation of a chosen completed category.
Cooper--Hogancamp Definition~5.3, p.~23, reverses both
gradings; their Remark~5.7, pp.~24--25, explicitly distinguishes
the bounded duality from the one-sided setting.
Thus the construction supplies no semisimple finite-Hom
ordinary self-adjoint realization excluded by (P-F52bm).

Finally let $\sum_\sigma a_\sigma Q_\sigma=0$ be the actual
source mixed relation, and take $\ell=11$, its maximum word
length.  Clearing the rescaling gives
\[
 \sum_\sigma a_\sigma F^{\ell-|\sigma|}
          [\mathscr Q_{\sigma_1}\cdots
             \mathscr Q_{\sigma_{|\sigma|}}]=0
                                                               \tag{P-F52bz}
\]
on the rational Euler subspace.  Every coefficient here is
an integral Laurent polynomial.  Splitting its positive and
negative terms gives two explicit finite sums of shifted
word functors.  The proof has not provided a natural homotopy
equivalence between those two functors, or compatible choices
for all mixed relations.  Equality of Euler classes does not
give such an equivalence: even a complex and its shift have
opposite Euler classes while their zero-differential direct
sum need not be contractible.  The known horizontal/vertical
fusion-category comparison in Cooper--Krushkal Theorem~6.6,
p.~23, is useful additional structure, but it does not state
this particular source relation or its required coherence.
Source canonical-basis multiplication, finite pure coefficient
models, and the proposed iterated-theta compatibility therefore
remain separate inputs.

\paragraph{The full tensor-image algebra, uniformly for the $\mathfrak{sl}_2$ family.}
The construction is not restricted to the four-dimensional module.
Fix integers $d,r\ge1$, put $X=V_d$, and let
\[
 B^{\mathrm{im}}_{d,r}
   =F_0\langle p^-_1,\ldots,p^-_{r-1}\rangle
      \ \subset\ A_{d,r}
      :=\operatorname{End}_{U_q\mathfrak{sl}_2}(V_d^{\otimes r})
                                                               \tag{P-F52ca}
\]
be the unital algebra of the actual tensor operators.
Here $p^-_i$ is the negative-braid spectral projector on
factors $i,i+1$; if $r=1$ there are no generators and the
algebra is $F_0$.
The notation $\mathrm{im}$ matters.  GCT VII, \S5,
printed pp.~13--14, first defines its algebra as a subalgebra
of a Hecke algebra.  Passage to the two-dimensional fundamental
tensor representation can have a kernel.  We make no general
claim that this Hecke subalgebra maps injectively to
(P-F52ca).

Let
$e=p_d\sqcup\cdots\sqcup p_d$ be the product of $r$
Jones--Wenzl idempotents in $\mathrm{TL}_{dr}(F_0)$.
There is a natural algebra isomorphism
\[
                e\,\mathrm{TL}_{dr}(F_0)\,e
                           \simeq A_{d,r}.
                                                               \tag{P-F52cb}
\]
Let $\mathscr C_{d,r}$ be the same completed category of
Temperley--Lieb tangle complexes, with
$\mathbf E=\mathbf P_d\sqcup\cdots\sqcup\mathbf P_d$
retained on both input and output boundaries.
Its completed Euler group has the multiplication induced
by gluing, and the cited Euler convention gives
\[
 K^{\mathrm{Eul}}(\mathscr C_{d,r})
       \otimes_{\widehat R}\mathbb C((q))
          \simeq A_{d,r}\otimes_{F_0}\mathbb C((q)).
                                                               \tag{P-F52cc}
\]
For each adjacent pair insert, with the boundary projectors
on both sides, the fixed complex
\[
        \mathbf N_d
          =\bigoplus_{\substack{0\le k\le d\\k\ {
m odd}}}
                      \mathbf P_{2d,\,2d-2k}.
                                                               \tag{P-F52cd}
\]
Left and right gluing with these complexes give actual
functors.  Their actions preserve the rational subspace
$B^{\mathrm{im}}_{d,r}$ in (P-F52cc), and on this subspace
are exactly left and right multiplication by $p^-_i$.
In particular, the left action is the faithful regular
representation of the full tensor-image algebra.
The Euler embedding preserves its multiplication and unit;
it is not merely an equality of the dimensions of two
representations.  We do not identify the whole Euler group
in (P-F52cc) with $B^{\mathrm{im}}_{d,r}$.

\emph{Proof.}
The fundamental Temperley--Lieb identification
$\mathrm{TL}_{dr}(F_0)=
\operatorname{End}_{U_q\mathfrak{sl}_2}(V_1^{\otimes dr})$
and the Jones--Wenzl image $p_dV_1^{\otimes d}=V_d$
give (P-F52cb).  Explicitly, restriction sends the corner
to endomorphisms of the image of $e$.  It is injective
because a corner endomorphism is zero on the complementary
kernel of $e$ and is determined by its image restriction.
It is surjective because an equivariant endomorphism of
that image extends to the ambient tensor space by composing
with the equivariant projection $e$.
These identifications use Cooper--Krushkal \S\S2.1--2.2,
printed pp.~4--5; the categorical compatibility with
composition is also stated in Cooper--Hogancamp
Theorem~3.1, p.~16.

The boundary Euler isomorphism is the general
Cooper--Krushkal Theorem~6.2 applied to the $2r$ labels
$(d,\ldots,d)$; bending the $r$ input bundles gives the
endomorphism version.  Gluing agrees with multiplication
by Lemma~2.10 and the composition statement just cited.
This proves (P-F52cc), with the same completion and rational
subspace convention as (P-F52bt).

For the projector in (P-F52cd), the tensor product is
multiplicity free:
$V_d\otimes V_d=\bigoplus_{k=0}^dV_{2d-2k}$.
The braid acts by a scalar of the form $\pm q^{a/2}$
on each summand, as in GCT VII, equation~(4), pp.~6--7.
Its sign is $(-1)^k$.  One can check the sign directly
at $q=1$: in the polynomial model for
$\operatorname{Sym}^d\mathbb C^2\otimes
 \operatorname{Sym}^d\mathbb C^2$, the highest vector
\[
 (x_1y_2-y_1x_2)^k x_1^{d-k}x_2^{d-k}
                                                               \tag{P-F52ce}
\]
belongs to the summand of highest weight $2d-2k$.
Interchanging the factors multiplies it by $(-1)^k$.
The braid specializes to that interchange, and a nonzero
scalar $\pm q^{a/2}$ has the same sign at $q=1$.
Thus (P-F52cd) has precisely the Euler class of the
negative spectral projector, with no choice of a
multiplicity-space splitting required.

Boundary preservation and idempotence follow from the
same higher-projector sliding and orthogonality used in
(P-F52bu).  Left or right gluing is a functor on the
completed tangle category.  Taking Euler classes gives
left or right multiplication in the corner algebra.
The word-generated rational subalgebra (P-F52ca) is
preserved under both actions.  Its left action is
faithful since an element annihilating its unit is zero.
All comparisons are algebraic over $F_0$ and remain
injective under Laurent expansion into $\mathbb C((q))$.
This proves the uniform tensor-image statement.
\hfill$\square$

For $d=r=3$, the full concrete source algebra is
21-dimensional: GCT VII, \S7.1.1, printed p.~19,
identifies its actual generators with the presentation
whose basis is the empty word and the twenty alternating
words of lengths one through ten.  Consequently this
entire algebra, with its regular left and right actions,
is the preserved Euler subspace of the
$p_3^{\sqcup3}$ corner of $\mathrm{TL}_9$.
The convenient generators $Q_i$ and the prescribed
generators $FQ_i$ are obtained by the finite shift
multiplicities $f$ and $Ff$, respectively.
This strengthens (P-F52bw) from its particular simple
module to the full concrete source algebra.  It still
does not give natural homotopies for all source word
relations, identify a source canonical basis with
indecomposable objects, or give finite pure coefficient
objects with the required positive multiplication.

\paragraph{An arithmetic obstruction to individual split-block Euler objects.}
There is a precise additional limitation on the completed construction
above.  Work before its scalar extension to $\mathbb C((q))$ and use
ordinary Temperley--Lieb diagrams as ambient coordinates.  An actual
admissible complex of finite sums of grading shifts of these diagrams
has Euler coordinates in
\[
             \widehat R=\mathbb Z[[q]][q^{-1}]
                          \ \subset\ \mathbb Q((q)).
                                                               \tag{P-F52cf}
\]
This applies also when its boundary projectors are the fixed infinite
complexes used above.  For each diagram and each power of $q$, the
convergence convention makes the Euler coefficient a finite alternating
sum of integer multiplicities.  Complex scalars in a differential or
in a cobordism do not change these multiplicities.  Here an ``actual
complex'' means an object of the stated diagram-complex category; we
are not adjoining new abstract objects by scalar extension of its
Euler group or by an additional idempotent completion.

In the cubic source example, no nonzero such Euler class can lie
wholly in either one of the two conjugate $G$-eigenspaces with
eigenvalues $g_1,g_2$.  This applies both to the full corner algebra
with left multiplication by $G$, and to the four-boundary module
with the inserted $G$ action.  In particular no nonzero scalar
multiple, even by an arbitrary element of $\mathbb C((q))$, of
either source vector $u_1,u_2$ can be the Euler class of an actual
complex in the corresponding full corner.  The same statement
holds for their nonzero images under the scalar extension of
(P-F52bw) to $\mathbb C((q))$.  It remains true
after every finite power regrading $q=t^k$, $k\ge1$.

\emph{Proof.}
All the rational tensor actions here are defined over $\mathbb Q(q)$.
Indeed the ordinary diagram composition has coefficients in
$\mathbb Z[q,q^{-1}]$, the Jones--Wenzl projectors have rational
$\mathbb Q(q)$ coefficients by their recurrence, and the channel
projectors have their rational Temperley--Lieb classes.  In the
cubic case $Q_i/f$ is the sum of the channel-$0$ and channel-$4$
classes.  Thus $Q_i$, and hence
$G=f^2-(Q_1-Q_2)^2$, have rational diagram coefficients.
Insertion or left multiplication by $G$ is consequently represented
by a matrix over $\mathbb Q(q)$ on the finite-dimensional ambient
diagram space.  These statements do not require rational
coordinates in a chosen fusion basis.

The discriminant and eigenvalues already identified from the exact
source relation and cell coefficients are
\[
\begin{split}
 D&=q^{24}(1+q^8)^2
       (5+8q^4-4q^6+18q^8-4q^{10}+8q^{12}+5q^{16}),\\
 R&=-3q^{12}-4q^{16}+2q^{18}-10q^{20}
                     +2q^{22}-4q^{24}-3q^{28},\\
 g_{1,2}&=\frac{-R\mp\sqrt D}{2q^{20}}.
\end{split}
                                                               \tag{P-F52cg}
\]
These are the same quantities verified in (P-F52f1)--(P-F52f3),
from GCT VII, Figure~2, printed p.~31, and Figure~7, p.~36.
The leading term of $D$ at zero is $5q^{24}$.
If a Laurent series over $\mathbb Q$ squared to $D$, its leading
exponent would be $12$ and its leading rational coefficient would
have square $5$, which is impossible.  Equivalently the two
eigenvalues, expanded in $\mathbb C((q))$ with the positive choice
of the leading coefficient of $\sqrt D$, begin as
\[
          g_{1,2}=\frac{3\mp\sqrt5}{2}\,q^{-8}
                               +O(q^{-4}),
             \qquad g_i\notin\mathbb Q((q)).
                                                               \tag{P-F52ch}
\]

Suppose a nonzero actual Euler vector $v$ belongs to one of these
eigenspaces.  By (P-F52cf) its ambient coordinates are in
$\mathbb Q((q))$, as are the coordinates of $Gv$.
Choose a coordinate $j$ for which $v_j\ne0$.  The equality
$Gv=g_i v$ implies
\[
                         g_i=(Gv)_j/v_j\in\mathbb Q((q)),
                                                               \tag{P-F52ci}
\]
contradicting (P-F52ch).  The earlier complete relation certificate
gives $Gu_i=g_i u_i$ and $u_i\ne0$; the scalar-extended map
(P-F52bw) is injective and intertwines $G$.  An arbitrary nonzero scalar multiple stays
in the same eigenspace, so the same argument excludes all the
stated scalar multiples without any restriction on their scalars.
After $q=t^k$, the leading term of $D(t^k)$ is $5t^{24k}$ and the
leading coefficients of $g_i(t^k)$ are still
$(3\pm\sqrt5)/2$.  Repeating the coordinate argument over
$\mathbb Q((t))$ proves the final assertion.
\hfill$\square$

Thus the scalar-extended Euler construction realizes the actual
algebra and its rational module, while an individual nonzero vector
in either split block cannot be represented by an actual integral
Euler object of this particular category.  Rational combinations
joining the conjugate blocks are not excluded; they are exactly
the type used in the descended source lattice.  This is a failure
of a specific object-by-object identification with the split source
vectors, not an obstruction to all geometric or cohomological
models.  In particular trace functions, different coefficient
structures, and different categories are not covered by this
argument.  Nor does it settle the mixed natural-homotopy relations
or the source's proposed canonical-basis products.
